% Kreisbogen-Überblendung an einem Stützpunkt eines stückweise linearen Gelenkraumpfads.
% Eigene Darstellung in Anlehnung an Kunz2012 (Fig. 1, S. 1; Belege Kunz2012.02, .03).
% Tangentenpunkte im Abstand 1,3 vom Stützpunkt, Mittelpunkt berechnet.
% Stil: ba-grafik.tex. Einbinden: % Kreisbogen-Überblendung an einem Stützpunkt eines stückweise linearen Gelenkraumpfads.
% Eigene Darstellung in Anlehnung an Kunz2012 (Fig. 1, S. 1; Belege Kunz2012.02, .03).
% Tangentenpunkte im Abstand 1,3 vom Stützpunkt, Mittelpunkt berechnet.
% Stil: ba-grafik.tex. Einbinden: % Kreisbogen-Überblendung an einem Stützpunkt eines stückweise linearen Gelenkraumpfads.
% Eigene Darstellung in Anlehnung an Kunz2012 (Fig. 1, S. 1; Belege Kunz2012.02, .03).
% Tangentenpunkte im Abstand 1,3 vom Stützpunkt, Mittelpunkt berechnet.
% Stil: ba-grafik.tex. Einbinden: % Kreisbogen-Überblendung an einem Stützpunkt eines stückweise linearen Gelenkraumpfads.
% Eigene Darstellung in Anlehnung an Kunz2012 (Fig. 1, S. 1; Belege Kunz2012.02, .03).
% Tangentenpunkte im Abstand 1,3 vom Stützpunkt, Mittelpunkt berechnet.
% Stil: ba-grafik.tex. Einbinden: \input{Bilder/tikz/kreisbogen-ueberblendung} in einer figure-Umgebung.
\begin{tikzpicture}[ba]
  \coordinate (A) at (0,0);
  \coordinate (B) at (4,2.2);
  \coordinate (C) at (7.5,0);
  \coordinate (T1) at (2.861,1.574);
  \coordinate (T2) at (5.101,1.508);
  \coordinate (M) at (3.925,-0.362);
  % ursprünglicher Pfad (Knick)
  \draw[ba hilfslinie] (T1) -- (B) -- (T2);
  % überblendeter Pfad
  \draw[ba glied, line width=1.6pt] (A) -- (T1) arc[start angle=118.8, end angle=57.8, radius=2.209] -- (C);
  % Radius
  \draw[ba duenn, baGrau] (M) -- (T2);
  \fill[baGrau] (M) circle (0.7mm);
  \node[anchor=north, text=baGrau] at (M) {Mittelpunkt};
  % Punkte
  \foreach \p/\l/\pos in {A/{q_{i-1}}/below, B/{q_i}/above, C/{q_{i+1}}/below}{
    \node[ba gelenk] at (\p) {};
    \node[\pos=1.5mm] at (\p) {$\l$};
  }
  \node[ba endeffektor, minimum size=2.4mm] at (T1) {};
  \node[ba endeffektor, minimum size=2.4mm] at (T2) {};
  \node[anchor=south east, align=right] at ($(T1)+(-0.2,0.15)$) {Beginn\\des Bogens};
  \node[anchor=south west, align=left] at ($(T2)+(0.2,0.15)$) {Ende\\des Bogens};
  % Legende
  \draw[ba glied, line width=1.6pt] (0,-1.3) -- ++(0.8,0); \node[anchor=west] at (0.9,-1.3) {überblendeter Pfad};
  \draw[ba hilfslinie] (4.6,-1.3) -- ++(0.8,0); \node[anchor=west] at (5.5,-1.3) {Pfad mit Knick};
\end{tikzpicture}
 in einer figure-Umgebung.
\begin{tikzpicture}[ba]
  \coordinate (A) at (0,0);
  \coordinate (B) at (4,2.2);
  \coordinate (C) at (7.5,0);
  \coordinate (T1) at (2.861,1.574);
  \coordinate (T2) at (5.101,1.508);
  \coordinate (M) at (3.925,-0.362);
  % ursprünglicher Pfad (Knick)
  \draw[ba hilfslinie] (T1) -- (B) -- (T2);
  % überblendeter Pfad
  \draw[ba glied, line width=1.6pt] (A) -- (T1) arc[start angle=118.8, end angle=57.8, radius=2.209] -- (C);
  % Radius
  \draw[ba duenn, baGrau] (M) -- (T2);
  \fill[baGrau] (M) circle (0.7mm);
  \node[anchor=north, text=baGrau] at (M) {Mittelpunkt};
  % Punkte
  \foreach \p/\l/\pos in {A/{q_{i-1}}/below, B/{q_i}/above, C/{q_{i+1}}/below}{
    \node[ba gelenk] at (\p) {};
    \node[\pos=1.5mm] at (\p) {$\l$};
  }
  \node[ba endeffektor, minimum size=2.4mm] at (T1) {};
  \node[ba endeffektor, minimum size=2.4mm] at (T2) {};
  \node[anchor=south east, align=right] at ($(T1)+(-0.2,0.15)$) {Beginn\\des Bogens};
  \node[anchor=south west, align=left] at ($(T2)+(0.2,0.15)$) {Ende\\des Bogens};
  % Legende
  \draw[ba glied, line width=1.6pt] (0,-1.3) -- ++(0.8,0); \node[anchor=west] at (0.9,-1.3) {überblendeter Pfad};
  \draw[ba hilfslinie] (4.6,-1.3) -- ++(0.8,0); \node[anchor=west] at (5.5,-1.3) {Pfad mit Knick};
\end{tikzpicture}
 in einer figure-Umgebung.
\begin{tikzpicture}[ba]
  \coordinate (A) at (0,0);
  \coordinate (B) at (4,2.2);
  \coordinate (C) at (7.5,0);
  \coordinate (T1) at (2.861,1.574);
  \coordinate (T2) at (5.101,1.508);
  \coordinate (M) at (3.925,-0.362);
  % ursprünglicher Pfad (Knick)
  \draw[ba hilfslinie] (T1) -- (B) -- (T2);
  % überblendeter Pfad
  \draw[ba glied, line width=1.6pt] (A) -- (T1) arc[start angle=118.8, end angle=57.8, radius=2.209] -- (C);
  % Radius
  \draw[ba duenn, baGrau] (M) -- (T2);
  \fill[baGrau] (M) circle (0.7mm);
  \node[anchor=north, text=baGrau] at (M) {Mittelpunkt};
  % Punkte
  \foreach \p/\l/\pos in {A/{q_{i-1}}/below, B/{q_i}/above, C/{q_{i+1}}/below}{
    \node[ba gelenk] at (\p) {};
    \node[\pos=1.5mm] at (\p) {$\l$};
  }
  \node[ba endeffektor, minimum size=2.4mm] at (T1) {};
  \node[ba endeffektor, minimum size=2.4mm] at (T2) {};
  \node[anchor=south east, align=right] at ($(T1)+(-0.2,0.15)$) {Beginn\\des Bogens};
  \node[anchor=south west, align=left] at ($(T2)+(0.2,0.15)$) {Ende\\des Bogens};
  % Legende
  \draw[ba glied, line width=1.6pt] (0,-1.3) -- ++(0.8,0); \node[anchor=west] at (0.9,-1.3) {überblendeter Pfad};
  \draw[ba hilfslinie] (4.6,-1.3) -- ++(0.8,0); \node[anchor=west] at (5.5,-1.3) {Pfad mit Knick};
\end{tikzpicture}
 in einer figure-Umgebung.
\begin{tikzpicture}[ba]
  \coordinate (A) at (0,0);
  \coordinate (B) at (4,2.2);
  \coordinate (C) at (7.5,0);
  \coordinate (T1) at (2.861,1.574);
  \coordinate (T2) at (5.101,1.508);
  \coordinate (M) at (3.925,-0.362);
  % ursprünglicher Pfad (Knick)
  \draw[ba hilfslinie] (T1) -- (B) -- (T2);
  % überblendeter Pfad
  \draw[ba glied, line width=1.6pt] (A) -- (T1) arc[start angle=118.8, end angle=57.8, radius=2.209] -- (C);
  % Radius
  \draw[ba duenn, baGrau] (M) -- (T2);
  \fill[baGrau] (M) circle (0.7mm);
  \node[anchor=north, text=baGrau] at (M) {Mittelpunkt};
  % Punkte
  \foreach \p/\l/\pos in {A/{q_{i-1}}/below, B/{q_i}/above, C/{q_{i+1}}/below}{
    \node[ba gelenk] at (\p) {};
    \node[\pos=1.5mm] at (\p) {$\l$};
  }
  \node[ba endeffektor, minimum size=2.4mm] at (T1) {};
  \node[ba endeffektor, minimum size=2.4mm] at (T2) {};
  \node[anchor=south east, align=right] at ($(T1)+(-0.2,0.15)$) {Beginn\\des Bogens};
  \node[anchor=south west, align=left] at ($(T2)+(0.2,0.15)$) {Ende\\des Bogens};
  % Legende
  \draw[ba glied, line width=1.6pt] (0,-1.3) -- ++(0.8,0); \node[anchor=west] at (0.9,-1.3) {überblendeter Pfad};
  \draw[ba hilfslinie] (4.6,-1.3) -- ++(0.8,0); \node[anchor=west] at (5.5,-1.3) {Pfad mit Knick};
\end{tikzpicture}
